\documentclass[12pt,a4paper]{article}
\usepackage{fontawesome5}
\providecommand{\faPencilAlt}{\faIcon{pencil-alt}}
\usepackage{styles/couleurs}
\usepackage{styles/mise_en_page}
\usepackage{styles/boites}
\usepackage{styles/maths}
\usepackage{styles/figures}
\usepackage{enumitem}
\usepackage{array}
\usepackage{tabularx}
\usepackage{amssymb}
\usepackage{pdflscape}
\usepackage{tikz}
\usepackage{xparse}
\setlist[itemize]{leftmargin=1.4em,itemsep=0.15em,topsep=0.2em}
\setlist[enumerate]{leftmargin=*,itemsep=0.2em,topsep=0.2em}
\newcommand{\lignepointillee}{\noindent\dotfill\par\vspace{0.48em}}
\newcommand{\zonelignes}[1]{%
\begin{tcolorbox}[enhanced,colback=white,colframe=NeutreC!40,boxrule=0.5pt,
arc=1.5mm,left=3mm,right=3mm,top=2mm,bottom=1mm,before skip=0pt,after skip=0pt]
\foreach \i in {1,...,#1}{\lignepointillee}
\end{tcolorbox}}
\newcounter{exohc}
\NewDocumentCommand{\exohc}{m +m}{%
\refstepcounter{exohc}\par\medskip
\noindent
\begin{minipage}[t]{0.47\linewidth}
\vspace{0pt}
{\sffamily\bfseries\color{TitreC}$\square$ Exercice \theexohc}
\vspace{0.2em}
#2
\end{minipage}
\hfill
\begin{minipage}[t]{0.50\linewidth}
\vspace{0pt}
\zonelignes{#1}
\end{minipage}
\par\medskip}
\begin{document}
\begin{landscape}

\begin{center}
{\sffamily\LARGE\bfseries\color{TitreC}Exercices hors classe}\\[0.2em]
{\sffamily\large\bfseries\color{ActionC}Calcul littéral}
\end{center}

\remarque{Cette fiche est donnée à l'avance.
Ne faire que les exercices indiqués par le professeur.}

\section*{A. Écrire et calculer une expression}

\exohc{5}{
Traduis par une expression littérale.
\begin{enumerate}[label=\alph*)]
  \item Le triple de \(x\), augmenté de \(7\).
  \item Le produit de \(5\) par la somme de \(x\) et de \(4\).
  \item La différence entre le carré de \(x\) et \(9\).
\end{enumerate}
}

\exohc{6}{
Un programme demande de choisir un nombre, de le multiplier par \(4\),
d'ajouter \(3\), puis de doubler le résultat.
\begin{enumerate}[label=\alph*)]
  \item Écris l'expression correspondant au nombre \(x\).
  \item Calcule pour \(x=5\).
  \item Calcule pour \(x=-2\).
\end{enumerate}
}

\exohc{5}{
Pour \(x=-3\), calcule :
\[
A=4x+7
\qquad
B=x^2-5
\qquad
C=2(x+6).
\]
}

\section*{B. Simplifier l'écriture}

\exohc{5}{
Simplifie.
\[
4\times x
\qquad
x\times7
\qquad
x\times x
\]
\[
1\times x+5
\qquad
3\times(x+2)
\qquad
2\times x\times x.
\]
}

\exohc{5}{
Explique pourquoi \(3+x\) et \(3x\) ne désignent pas la même expression.
Choisis une valeur de \(x\) pour le vérifier.
}

\section*{C. Réduire}

\exohc{6}{
Réduis.
\[
A=3x+5x
\qquad
B=x+x+x+4
\]
\[
C=7x+3-2x+5
\qquad
D=-4x+9+6x-2.
\]
}

\exohc{6}{
Réduis.
\[
E=5x^2+2x-3x^2+7x
\]
\[
F=-3x^2+4x+8x^2-9x+2.
\]
}

\exohc{5}{
Deux élèves écrivent :
\[
3x+5=8x
\qquad
2x+3x^2=5x^3.
\]
Explique et corrige les deux erreurs.
}

\section*{D. Développer puis réduire}

\exohc{6}{
Développe.
\[
A=5(x+2)
\qquad
B=3(2x+7)
\]
\[
C=-2(x+4)
\qquad
D=4(3x-5).
\]
}

\exohc{7}{
Développe puis réduis.
\[
E=3(x+4)+2x
\]
\[
F=5(2x-1)-3x
\]
\[
G=-2(x+3)+7x-4.
\]
}

\section*{E. Prouver ou réfuter}

\exohc{7}{
Montre que les expressions suivantes sont égales pour toute valeur de \(x\) :
\[
4(x+3)+2x
\qquad
6x+12.
\]
}

\exohc{6}{
L'égalité suivante est-elle vraie pour toute valeur de \(x\) ?
\[
3(x+2)=3x+2.
\]
Justifie par un contre-exemple.
}

\end{landscape}
\end{document}
